\section{Further examples: where the standard laws come from}
\label{sec:distribution-provenance}

The distributions introduced earlier are descriptions of experiments,
not just lists of formulas. The table makes their origins explicit.
The negative-binomial, gamma, and order-statistic beta examples that
follow are enrichment; the continuous parts use the calculus refresher.

\begin{center}\small
\begin{tabularx}{\textwidth}{p{.23\textwidth}X}
\toprule
Law & Generating experiment and assumptions \\
\midrule
Discrete uniform & One uniformly chosen label in a finite list. \\
Bernoulli & Record any event of probability $p$ as zero or one. \\
Binomial & Count successes in a fixed number of independent trials with common $p$. \\
Hypergeometric & Count marked objects in a uniform sample without replacement. \\
Geometric & Count independent trials through the first success, with fixed $p>0$. \\
Poisson & Count arrivals in a fixed interval of a constant-rate Poisson model; also a limit of many rare independent trials. \\
Continuous uniform & Choose a point with constant density on an interval. \\
Exponential & Wait for the first arrival in a constant-rate Poisson model. \\
Normal & An exact Gaussian model when assumed; also an approximation to standardized sums under the central limit theorem. \\
\bottomrule
\end{tabularx}
\end{center}

\subsection{Negative binomial: wait for the \texorpdfstring{$r$}{r}th success}
\label{sec:negative-binomial}

Use independent Bernoulli trials with success probability $0<p\le1$.
For an integer $r\ge1$, let $V_r$ count trials through the $r$th success.
For $0<p<1$ and $k=r,r+1,\ldots$,
\[
\boxed{\P(V_r=k)=\binom{k-1}{r-1}p^r(1-p)^{k-r}.}
\]
The last trial must succeed, and the first $k-1$ trials must contain
exactly $r-1$ successes. This proves the mass formula from its experiment.
Alternatively, $V_r$ is the sum of $r$ independent geometric waits:
after each success the unused trials are a fresh independent sequence.
Each wait is finite almost surely, so the displayed masses exhaust all
outcomes and sum to one. The geometric moments give
\[
\E[V_r]=r/p,\qquad \Var(V_r)=r(1-p)/p^2.
\]
At $p=1$, $V_r=r$ with certainty. For $r=1$ this is the geometric law.
Some authors instead call the \emph{failure count} $V_r-r$ negative
binomial; its support starts at zero and its mean is $r(1-p)/p$.

\begin{example}[Three successful attempts]
With independent success probability $1/2$, the number of attempts
through success number three has mean $6$, variance $6$, and
$\P(V_3=5)=\binom42(1/2)^5=3/16$.
The experiment fixes the number of successes, whereas a binomial
experiment fixes the number of trials.
\end{example}

\subsection{Uniform and exponential waits come from different clocks}

If a bus departs every $b$ minutes and one's arrival phase is uniform
on $[0,b]$, the wait is Uniform$(0,b)$, with mean $b/2$ and variance
$b^2/12$. Uniform phase is a modeling assumption; a timetable alone
does not make a passenger's arrival phase uniform.

By contrast, suppose counts in an interval of length $t$ are
Poisson$(\lambda t)$, with independent counts in disjoint intervals.
The wait $E$ for the next arrival exceeds $t$ exactly when there are
zero arrivals, so
\[
\P(E>t)=e^{-\lambda t},\quad
f_E(t)=\lambda e^{-\lambda t},\quad t\ge0.
\]
This is the exponential law with \emph{rate} $\lambda$ and mean
$1/\lambda$. Division of survival probabilities proves memorylessness:
$\P(E>s+t\mid E>s)=e^{-\lambda t}$.
For calls modeled at rate $3$ per hour, 20 minutes has mean count one,
probability $1-2e^{-1}$ of at least two calls, and probability $e^{-1}$
of waiting more than 20 minutes for the next call.

There is also a precise discrete-to-continuous connection. If
$G_n\sim\text{Geometric}(\lambda/n)$, then for fixed $t\ge0$ and
integers $n>\lambda$,
\[
\P(G_n/n>t)=(1-\lambda/n)^{\lfloor nt\rfloor}
\longrightarrow e^{-\lambda t}.
\]
Thus an exponential wait is a limit of rescaled rare-success waits.
The equality before the limit concerns an integer-trial experiment;
the exponential survival function is its limiting continuous model.

\subsection{Gamma: add exponential waiting times}
\label{sec:gamma-waits}

Let $E_1,E_2,\ldots$ be independent exponential variables of rate
$\lambda>0$, and put $S_r=E_1+\cdots+E_r$. Then, for integer $r\ge1$,
\[
\boxed{f_{S_r}(t)=\frac{\lambda^r t^{r-1}e^{-\lambda t}}{(r-1)!},
\qquad t>0.}
\]
This is Gamma$(r,\text{rate }\lambda)$, also called an Erlang law.
To derive it, start with $r=1$. Conditioning on the previous subtotal
and integrating the independent next-wait density gives
\begin{align*}
f_{S_{r+1}}(t)
&=\int_0^t f_{S_r}(s)\lambda e^{-\lambda(t-s)}\,ds\\
&=\frac{\lambda^{r+1}e^{-\lambda t}}{(r-1)!}
\int_0^t s^{r-1}\,ds
=\frac{\lambda^{r+1}t^re^{-\lambda t}}{r!}.
\end{align*}
This single-integral formula is called \emph{convolution}; it sums
the possibilities for the first subtotal. Integration by parts gives
$\int_0^\infty t^j e^{-\lambda t}\,dt=j!/\lambda^{j+1}$ for integers
$j\ge0$, verifying normalization and the moments
\[
\E[S_r]=r/\lambda,\qquad \Var(S_r)=r/\lambda^2.
\]
They also follow by adding the independent exponential moments.
General positive noninteger gamma shapes exist too; the arrival-count
interpretation here uses integer $r$.

Integrating the density, or differentiating the expression below and
checking its value at zero, gives
\[
\P(S_r>t)=e^{-\lambda t}\sum_{j=0}^{r-1}\frac{(\lambda t)^j}{j!}.
\]
If arrivals occur at times $S_1,S_2,\ldots$, their count $N(t)$ obeys
$\{N(t)\ge r\}=\{S_r\le t\}$. Taking differences of these probabilities
gives $\P(N(t)=k)=e^{-\lambda t}(\lambda t)^k/k!$ for $k\ge0$.
Thus exponential gaps construct the Poisson count model. The gamma
tail tends to one as $r\to\infty$ at fixed $t$, so only finitely many
arrivals occur in a bounded interval, almost surely.

For rate $2$ per minute, the wait for the third arrival has mean
$3/2$ minutes, variance $3/4$ minutes squared, and
$\P(S_3>1)=e^{-2}(1+2+2^2/2)=5e^{-2}$.

\subsection{Beta: the location of an ordered uniform sample}
\label{sec:beta-order-statistics}

Choose $m$ independent Uniform$(0,1)$ points and let $U_{(k)}$ be the
$k$th smallest, where $1\le k\le m$. For $0\le x\le1$, at least $k$
points lie at or below $x$ exactly when $U_{(k)}\le x$. Hence
\[
\P(U_{(k)}\le x)=\sum_{j=k}^m\binom mjx^j(1-x)^{m-j}.
\]
Differentiate this finite sum. The identities
$j\binom mj=m\binom{m-1}{j-1}$ and
$(m-j)\binom mj=m\binom{m-1}j$ cancel adjacent terms, leaving
\[
\boxed{f_{U_{(k)}}(x)=\frac{m!}{(k-1)!(m-k)!}
x^{k-1}(1-x)^{m-k},\qquad 0<x<1.}
\]
This is Beta$(k,m+1-k)$. Its normalization follows from the CDF;
the integer beta integral in Appendix~\ref{app:calculus} gives
\[
\E[U_{(k)}]=\frac{k}{m+1},\qquad
\Var(U_{(k)})=\frac{k(m+1-k)}{(m+1)^2(m+2)}.
\]
For example, the middle of three uniform points is Beta$(2,2)$, with
density $6x(1-x)$, mean $1/2$, and variance $1/20$. Its chance of
lying between $1/4$ and $3/4$ is $11/16$, using the CDF $3x^2-2x^3$.
This sampling origin of beta laws is different from their use as priors
in Chapter 9; neither interpretation is forced by the name alone.

\subsection{Normal: a model and a limiting shape}

For independent copies $X_i$ of a variable with finite mean $\mu$ and
finite nonzero variance $\sigma^2$, the central limit theorem in
Chapter 7 states that
\[
\frac{X_1+\cdots+X_n-n\mu}{\sigma\sqrt n}
\quad\text{converges in distribution to }\Normal(0,1).
\]
This explains the normal approximation to many accumulated small
contributions. It requires assumptions; a common random shock or an
infinite variance can invalidate this particular theorem.

For 100 independent fair tosses, the head count is exactly
$\Bin(100,1/2)$, with mean $50$ and standard deviation $5$.
The exact probability of 45 through 55 heads is
$2^{-100}\sum_{j=45}^{55}\binom{100}j$.
A continuity-corrected normal approximation is
$\Phi(1.1)-\Phi(-1.1)$: integrate from $44.5$ to $55.5$ before
standardizing. The binomial law is exact under the toss assumptions;
the normal answer is an approximation. For continuous measurement errors,
normality can instead be an explicit model assumption to assess from the
mechanism and data, rather than a conclusion from a bell-shaped histogram.
